\section{Vectors from Coordinate Displacements}

In the previous chapter a point of the Euclidean plane was identified with an
ordered pair of real numbers. We now use that identification immediately. The
aim is not to collect formulas about vectors, but to let the formulas emerge
from geometric questions and coordinate calculations.

Choose two points
\[
A=(1,2),
\qquad
B=(4,6).
\]
To pass from $A$ to $B$, the first coordinate changes by $3$ and the second by
$4$. The pair
\[
(3,4)
\]
records the complete displacement.

\subsection{A Vector Records a Change of Coordinates}

For arbitrary points
\[
A=(x_1,y_1),
\qquad
B=(x_2,y_2),
\]
the displacement from $A$ to $B$ is
\[
\overrightarrow{AB}
=
(x_2-x_1,\,y_2-y_1).
\]

\begin{definitionbox}
A vector in the Cartesian plane is an ordered pair
\[
\mathbf v=(v_1,v_2)
\]
interpreted as a coordinate displacement: $v_1$ units horizontally and $v_2$
units vertically.
\end{definitionbox}

This makes the phrase ``the same vector starting elsewhere'' precise. If
$P=(p_1,p_2)$ and
\[
Q=(p_1+v_1,\,p_2+v_2),
\]
then
\[
\overrightarrow{PQ}=(v_1,v_2)=\mathbf v.
\]
The same two coordinate changes have simply been applied at another point.

\begin{center}
\begin{tikzpicture}[scale=0.85]
    \draw[axis] (-0.5,0) -- (7.0,0) node[right] {$x$};
    \draw[axis] (0,-0.5) -- (0,5.0) node[above] {$y$};
    \coordinate (A) at (1,1);
    \coordinate (B) at (4,3);
    \coordinate (P) at (2.5,2.2);
    \coordinate (Q) at (5.5,4.2);
    \draw[subset,->] (A) -- (B) node[midway,below right] {$(3,2)$};
    \draw[subset,->] (P) -- (Q) node[midway,above left] {$(3,2)$};
    \fill[geometricSubsetColor] (A) circle (2pt) node[below left] {$A$};
    \fill[geometricSubsetColor] (B) circle (2pt) node[below right] {$B$};
    \fill[geometricSubsetColor] (P) circle (2pt) node[above left] {$P$};
    \fill[geometricSubsetColor] (Q) circle (2pt) node[above right] {$Q$};
\end{tikzpicture}
\end{center}

The zero displacement is
\[
\mathbf0=(0,0).
\]
Reversing $\mathbf v=(v_1,v_2)$ gives
\[
-\mathbf v=(-v_1,-v_2).
\]

\subsection{Two Successive Displacements Force Vector Addition}

Start at an arbitrary point $(x,y)$. First apply
\[
\mathbf u=(u_1,u_2),
\]
and then
\[
\mathbf v=(v_1,v_2).
\]
The successive endpoints are
\[
(x,y)
\longmapsto
(x+u_1,y+u_2)
\longmapsto
(x+u_1+v_1,y+u_2+v_2).
\]
Therefore the single displacement with the same effect must be
\[
\mathbf u+\mathbf v
=
(u_1+v_1,\,u_2+v_2).
\]
The rule is forced by the experiment; it is not introduced as a convention to
memorise.

Because addition of real numbers is commutative, applying $\mathbf u$ and then
$\mathbf v$ reaches the same endpoint as applying $\mathbf v$ and then
$\mathbf u$. This produces the parallelogram rule.

\begin{center}
\begin{tikzpicture}[scale=0.95]
    \coordinate (O) at (0,0);
    \coordinate (U) at (3.2,0.9);
    \coordinate (V) at (1.0,2.5);
    \coordinate (S) at (4.2,3.4);
    \draw[subset,->] (O) -- (U) node[midway,below] {$\mathbf u$};
    \draw[subset,->] (O) -- (V) node[midway,left] {$\mathbf v$};
    \draw[helper] (U) -- (S);
    \draw[helper] (V) -- (S);
    \draw[very thick,->] (O) -- (S) node[midway,above left] {$\mathbf u+\mathbf v$};
\end{tikzpicture}
\end{center}

Subtraction is the displacement that undoes another displacement:
\[
\mathbf u-\mathbf v
=
(u_1-v_1,\,u_2-v_2).
\]
In particular,
\[
\overrightarrow{AB}=B-A.
\]

\subsection{Repeating and Reversing a Displacement}

Repeating $\mathbf v=(v_1,v_2)$ twice gives $(2v_1,2v_2)$, while reversing it
gives $(-v_1,-v_2)$. The natural extension is
\[
\lambda\mathbf v=(\lambda v_1,\lambda v_2),
\qquad \lambda\in\mathbb R.
\]
For positive $\lambda$ the orientation is unchanged; for negative $\lambda$ it
is reversed.

\subsection{Length and Inclination}

Draw $\mathbf v=(v_1,v_2)$ from the origin. Its endpoint and the coordinate axes
form a right triangle. The Pythagorean theorem gives
\[
\|\mathbf v\|^2=v_1^2+v_2^2,
\qquad
\|\mathbf v\|=\sqrt{v_1^2+v_2^2}.
\]
Consequently,
\[
d(A,B)=\|B-A\|.
\]

\begin{center}
\begin{tikzpicture}[scale=0.9]
    \coordinate (O) at (0,0);
    \coordinate (P) at (4,2.7);
    \coordinate (R) at (4,0);
    \draw[subset,->] (O) -- (P) node[midway,above left] {$\mathbf v$};
    \draw[helper] (O) -- (R) node[midway,below] {$|v_1|$};
    \draw[helper] (R) -- (P) node[midway,right] {$|v_2|$};
    \draw (R) ++(-0.28,0) -- ++(0,0.28) -- ++(0.28,0);
\end{tikzpicture}
\end{center}

Let $\alpha$ be the angle of inclination of a nonzero vector $\mathbf v$ from
the positive horizontal direction. The coordinate projections of $\mathbf v$
are
\[
v_1=\|\mathbf v\|\cos\alpha,
\qquad
v_2=\|\mathbf v\|\sin\alpha.
\]
Thus
\[
\mathbf v
=
\|\mathbf v\|(\cos\alpha,\sin\alpha).
\]
This representation applies to every nonzero vector; no vector has been chosen
as a special coordinate axis.

\subsection{One Angle Calculation, Two Geometric Tests}

Take two arbitrary nonzero vectors
\[
\mathbf u=(u_1,u_2),
\qquad
\mathbf v=(v_1,v_2).
\]
Let their angles of inclination be $\alpha$ and $\beta$. Then
\[
\mathbf u
=
\|\mathbf u\|(\cos\alpha,\sin\alpha),
\]
\[
\mathbf v
=
\|\mathbf v\|(\cos\beta,\sin\beta).
\]
The difference $\beta-\alpha$ measures the directed turn from $\mathbf u$ to
$\mathbf v$. The ordinary angle between the vectors will be denoted by
\[
\theta=|\beta-\alpha|
\]
when the smaller angle is intended.

\begin{center}
\begin{tikzpicture}[scale=0.9]
    \coordinate (O) at (0,0);
    \coordinate (U) at (4.2,1.6);
    \coordinate (V) at (1.6,3.6);
    \draw[axis] (-0.4,0) -- (5.2,0) node[right] {$x$};
    \draw[subset,->] (O) -- (U) node[above right] {$\mathbf u$};
    \draw[subset,->] (O) -- (V) node[above] {$\mathbf v$};
    \draw (0.8,0) arc[start angle=0,end angle=20.9,radius=0.8];
    \node at (0.97,0.18) {$\alpha$};
    \draw (0.66,0.25) arc[start angle=20.9,end angle=66.0,radius=0.71];
    \node at (0.61,0.62) {$\theta$};
    \draw (1.35,0) arc[start angle=0,end angle=66.0,radius=1.35];
    \node at (1.02,1.08) {$\beta$};
\end{tikzpicture}
\end{center}

\subsubsection{Parallel directions arise from the sine of the angle difference}

Two nonzero vectors determine the same unoriented direction precisely when the
turn from one to the other is an integer multiple of a half-turn:
\[
\beta-\alpha\equiv0\pmod{\pi}.
\]
Equivalently,
\[
\sin(\beta-\alpha)=0.
\]
Using the sine-difference identity,
\[
\sin(\beta-\alpha)
=
\sin\beta\cos\alpha-
\cos\beta\sin\alpha.
\]
Now substitute
\[
\cos\alpha=\frac{u_1}{\|\mathbf u\|},
\qquad
\sin\alpha=\frac{u_2}{\|\mathbf u\|},
\]
\[
\cos\beta=\frac{v_1}{\|\mathbf v\|},
\qquad
\sin\beta=\frac{v_2}{\|\mathbf v\|}.
\]
We obtain
\[
\begin{aligned}
\sin(\beta-\alpha)
&=
\frac{v_2}{\|\mathbf v\|}
\frac{u_1}{\|\mathbf u\|}
-
\frac{v_1}{\|\mathbf v\|}
\frac{u_2}{\|\mathbf u\|}\\
&=
\frac{u_1v_2-u_2v_1}
{\|\mathbf u\|\,\|\mathbf v\|}.
\end{aligned}
\]
The denominator is nonzero. Therefore
\[
\sin(\beta-\alpha)=0
\iff
u_1v_2-u_2v_1=0.
\]
We have discovered the coordinate test
\[
\mathbf u\parallel\mathbf v
\iff
u_1v_2-u_2v_1=0.
\]

The same condition implies that one vector is a scalar multiple of the other.
Indeed, if $v_1\neq0$, set $\lambda=u_1/v_1$. The equality
\[
u_1v_2-u_2v_1=0
\]
gives $u_2=\lambda v_2$, and hence
\[
\mathbf u=\lambda\mathbf v.
\]
If $v_1=0$, then $v_2\neq0$ and the same argument uses
$\lambda=u_2/v_2$. Thus
\[
\mathbf u\parallel\mathbf v
\iff
\mathbf u=\lambda\mathbf v
\quad\text{for some }\lambda\neq0.
\]
This scalar-multiple criterion is a conclusion of the angle calculation, not a
rule imposed in advance.

\subsubsection{The cosine of the same difference reveals perpendicularity}

The vectors are perpendicular precisely when their angle difference is a
quarter-turn modulo a half-turn:
\[
\beta-\alpha
\equiv
\frac{\pi}{2}
\pmod{\pi}.
\]
Equivalently,
\[
\cos(\beta-\alpha)=0.
\]
The cosine-difference identity gives
\[
\cos(\beta-\alpha)
=
\cos\beta\cos\alpha+
\sin\beta\sin\alpha.
\]
Substituting the coordinate ratios gives
\[
\begin{aligned}
\cos(\beta-\alpha)
&=
\frac{v_1}{\|\mathbf v\|}
\frac{u_1}{\|\mathbf u\|}
+
\frac{v_2}{\|\mathbf v\|}
\frac{u_2}{\|\mathbf u\|}\\
&=
\frac{u_1v_1+u_2v_2}
{\|\mathbf u\|\,\|\mathbf v\|}.
\end{aligned}
\]
Since the denominator is nonzero,
\[
\mathbf u\perp\mathbf v
\iff
u_1v_1+u_2v_2=0.
\]
The coordinate expression has emerged from the geometry. Only now do we give it
a name.

\begin{definitionbox}
For vectors in the plane, their \emph{dot product} is
\[
\mathbf u\cdot\mathbf v
=
u_1v_1+u_2v_2.
\]
In three-dimensional space,
\[
(u_1,u_2,u_3)\cdot(v_1,v_2,v_3)
=
u_1v_1+u_2v_2+u_3v_3.
\]
\end{definitionbox}

The preceding calculation proves
\[
\mathbf u\perp\mathbf v
\iff
\mathbf u\cdot\mathbf v=0.
\]
It proves more than the right-angle case. Because cosine is even and
$\theta=|\beta-\alpha|$,
\[
\cos(\beta-\alpha)=\cos\theta.
\]
Therefore
\[
\boxed{\mathbf u\cdot\mathbf v
=
\|\mathbf u\|\,\|\mathbf v\|\cos\theta}
\]
for every pair of nonzero vectors in the plane.

\begin{examplebox}
Take
\[
\mathbf u=(3,2),
\qquad
\mathbf v=(-2,3).
\]
Then
\[
\mathbf u\cdot\mathbf v
=3(-2)+2(3)=0.
\]
The coordinate calculation says that
\[
\cos\theta=0,
\]
so the angle between the vectors is a right angle.
\end{examplebox}

\subsection{The Cosine Theorem Emerges Immediately}

Draw $\mathbf u$ and $\mathbf v$ from a common point. The displacement from the
endpoint of $\mathbf v$ to the endpoint of $\mathbf u$ is
\[
\mathbf u-\mathbf v.
\]
Thus these three vectors determine a triangle with side lengths
\[
\|\mathbf u\|,
\qquad
\|\mathbf v\|,
\qquad
\|\mathbf u-\mathbf v\|.
\]
We calculate the square of the third side:
\[
\begin{aligned}
\|\mathbf u-\mathbf v\|^2
&=(\mathbf u-\mathbf v)\cdot(\mathbf u-\mathbf v)\\
&=\mathbf u\cdot\mathbf u
-\mathbf u\cdot\mathbf v
-\mathbf v\cdot\mathbf u
+\mathbf v\cdot\mathbf v\\
&=\|\mathbf u\|^2+
\|\mathbf v\|^2-
2\mathbf u\cdot\mathbf v\\
&=\|\mathbf u\|^2+
\|\mathbf v\|^2-
2\|\mathbf u\|\,\|\mathbf v\|\cos\theta.
\end{aligned}
\]
If
\[
a=\|\mathbf u\|,
\qquad
b=\|\mathbf v\|,
\qquad
c=\|\mathbf u-\mathbf v\|,
\]
then
\[
\boxed{c^2=a^2+b^2-2ab\cos\theta}.
\]

\begin{theorembox}
For a triangle with side lengths $a$, $b$, and $c$, where $\theta$ is the angle
between the sides of lengths $a$ and $b$,
\[
c^2=a^2+b^2-2ab\cos\theta.
\]
\end{theorembox}

The cosine theorem is not inserted as an external formula. It is the coordinate
expansion of the displacement between the endpoints of two arbitrary vectors.

\begin{center}
\begin{tikzpicture}[scale=0.9]
    \coordinate (O) at (0,0);
    \coordinate (U) at (4.4,1.5);
    \coordinate (V) at (1.3,3.5);
    \draw[subset,->] (O) -- (U) node[midway,below] {$\mathbf u$};
    \draw[subset,->] (O) -- (V) node[midway,left] {$\mathbf v$};
    \draw[very thick,->] (V) -- (U) node[midway,above right] {$\mathbf u-\mathbf v$};
    \draw (0.72,0.25) arc[start angle=19,end angle=69.6,radius=0.76];
    \node at (0.70,0.65) {$\theta$};
\end{tikzpicture}
\end{center}

\subsection{Projection Is Forced by the Right-Angle Condition}

Let $\mathbf v\neq\mathbf0$. We seek a multiple $\lambda\mathbf v$ such that
the remainder
\[
\mathbf u-\lambda\mathbf v
\]
is perpendicular to $\mathbf v$. The criterion just derived requires
\[
(\mathbf u-\lambda\mathbf v)\cdot\mathbf v=0.
\]
Expanding gives
\[
\mathbf u\cdot\mathbf v-
\lambda(\mathbf v\cdot\mathbf v)=0,
\]
so
\[
\lambda
=
\frac{\mathbf u\cdot\mathbf v}{\mathbf v\cdot\mathbf v}.
\]
Thus the parallel component is
\[
\operatorname{proj}_{\mathbf v}\mathbf u
=
\frac{\mathbf u\cdot\mathbf v}{\mathbf v\cdot\mathbf v}\,\mathbf v.
\]
The formula is the solution of a geometric requirement, not a rule to be
memorised without explanation.

\subsection{The Same Coordinate Ideas in Three Dimensions}

In $\mathbb R^3$, a vector is a displacement
\[
\mathbf v=(v_1,v_2,v_3),
\]
and addition and scaling remain componentwise. Applying the Pythagorean theorem
twice gives
\[
\|\mathbf v\|
=
\sqrt{v_1^2+v_2^2+v_3^2}.
\]

Two nonparallel vectors in space determine two independent directions. To
construct a vector perpendicular to both, solve the corresponding two dot-product
conditions. The resulting expression is
\[
\mathbf u\times\mathbf v
=
\bigl(
 u_2v_3-u_3v_2,
 u_3v_1-u_1v_3,
 u_1v_2-u_2v_1
\bigr).
\]
Direct substitution verifies
\[
(\mathbf u\times\mathbf v)\cdot\mathbf u=0,
\qquad
(\mathbf u\times\mathbf v)\cdot\mathbf v=0.
\]
Thus the cross product is introduced as a constructed solution to two explicit
perpendicularity conditions.

\begin{summarybox}
A vector was introduced as a coordinate displacement. Successive displacements
forced vector addition, and repetition forced scalar multiplication. Writing
arbitrary vectors through their lengths and inclination angles allowed the sine
of the angle difference to produce the parallelism test and the cosine of the
same difference to produce the dot product, perpendicularity, and the formula
\[
\mathbf u\cdot\mathbf v
=
\|\mathbf u\|\,\|\mathbf v\|\cos\theta.
\]
Expanding $\|\mathbf u-\mathbf v\|^2$ then produced the cosine theorem. Projection
and the cross product arose by solving explicit perpendicularity conditions.
\end{summarybox}
